% MR55--0C
% Kollaps/Synkope
% 14.0
% 2013-11-30

:ref:`m-kollaps`Taktik
    Ausführliche Untersuchung und symptomatische Behandlung 
\Ca{Verwechslung mit lebensbedrohlichen Differentialdiagnosen möglich!}}

\begin{itemize}
  -   Vitale Bedrohung einschätzen, auf **Differentialdiagnosen** untersuchen, z.\,B.:
  \begin{compactitem}
    -   Insult
    -   Herzrhytmusstörungen
    -   Akutes Koronarsyndrom (evtl. schmerzlos!)
    -   Krampfanfall
    -   Schädel-Hirn-Trauma (SHT)
    -   Zuckerstoffwechselstörung (Hyper-, Hypoglykämie)
    -   Exsikkose
    -   \ldots
  \end{compactitem}
  Es sind dabei alle zur Verfügung stehenden
  Möglichkeiten auszuschöpfen, z.\,B.:
  \begin{compactitem}
    -   Traumacheck (Sturz?),
    -   Neurocheck inkl. Blutzuckermessung,
    -   Temperatur,
    -   Ausführliche Anamnese bzw. Fremdanamnese 
    -   \ldots
  \end{compactitem}
  -   **Ursachenforschung** (Begleiterkrankungen (Infekt, \ldots), letzte Mahlzeit, Hitze, \ldots)
  -   Lagerung: Nach Ausschluss von Herzbeschwerden,
  Atemnot und relevanten Verletzungen: Beine hoch
  -   Transportentscheidung: Hospitalisierung zur Abklärung
  anstreben\newline
  Abteilung: Innere Medizin
\end{itemize}